% Original conceptual curves; not fitted to the lecture's experimental data.
\begin{tikzpicture}[>=Latex,
  axis/.style={->,thick},
  curve/.style={very thick},
  title/.style={font=\small\bfseries},
  note/.style={font=\scriptsize,align=center}
]
  \begin{scope}[xshift=0cm]
    \node[title] at (1.8,3.0) {$\alpha$ too small};
    \draw[axis] (0,0) -- (3.7,0) node[right] {iteration};
    \draw[axis] (0,0) -- (0,2.6) node[above] {$J$};
    \draw[curve,draw=noteBlue]
      plot[smooth] coordinates {(0.2,2.2)(0.8,1.85)(1.4,1.55)(2.0,1.30)(2.6,1.10)(3.3,0.93)};
    \node[note] at (1.8,-0.45) {stable but slow};
  \end{scope}

  \begin{scope}[xshift=4.5cm]
    \node[title] at (1.8,3.0) {stable $\alpha$};
    \draw[axis] (0,0) -- (3.7,0) node[right] {iteration};
    \draw[axis] (0,0) -- (0,2.6) node[above] {$J$};
    \draw[curve,draw=noteBlue]
      plot[smooth] coordinates {(0.2,2.2)(0.7,1.25)(1.2,0.68)(1.8,0.36)(2.5,0.21)(3.3,0.16)};
    \node[note] at (1.8,-0.45) {fast, monotone decrease};
  \end{scope}

  \begin{scope}[xshift=9.0cm]
    \node[title] at (1.8,3.0) {$\alpha$ too large};
    \draw[axis] (0,0) -- (3.7,0) node[right] {iteration};
    \draw[axis] (0,0) -- (0,2.6) node[above] {$J$};
    \draw[curve,draw=ntuRed]
      plot[smooth] coordinates {(0.2,2.1)(0.6,0.55)(1.0,2.25)(1.4,0.40)(1.8,2.35)(2.2,0.32)(2.6,2.45)(3.2,0.25)};
    \node[note] at (1.8,-0.45) {oscillation or divergence};
  \end{scope}
\end{tikzpicture}
